% Proposed insertion after the recovery section. NOT included by main.tex.
% These new results are not Lean-checked. Before inclusion, revise the abstract,
% formalization scope, and correspondence table rather than extending their
% kernel-verification claim silently. The geometric envelope-rate extension is
% proved separately in 03-nonconvex-recovery.md and is not included here.
\section{Stability}\label{sec:stability}

Write $M=\abs\Gs$. The equality argument has a quantitative form: the deficit
in $\cQ$ controls the support function of the cap, modulo horizontal
translation. We first give an explicit estimate on $\cT$, and then a
qualitative stability theorem for arbitrary moving sofas. The latter has no
injectivity hypothesis, but supplies no rate.

\begin{theorem}[quantitative cap rigidity]\label{thm:stability-cap}
Let $\xi=(K,B,D)\in\cT$, put $\Delta=h_K-h_{\KG}$ and $s=-\Delta(\pi)$.
Then
\begin{equation}\label{eq:stability-Q}
  \dH\bigl(K,\KG+(s,0)\bigr)
    \le 2\sec\varphi\,\sqrt{M-\cQ(\xi)}.
\end{equation}
In particular, every $K\in\Ki$ satisfies
\begin{equation}\label{eq:stability-area}
  \dH\bigl(K,\KG+(s,0)\bigr)
    \le 2\sec\varphi\,\sqrt{M-\cA(K)}.
\end{equation}
On the parameter interval $\varphi\in[0.039,0.04]$, the constant
$2\sec\varphi$ is less than $1001/500=2.002$.
\end{theorem}

\begin{lemma}[the pinned residual estimate]\label{lem:stability-residual}
Let $0<\varphi<\pi/4$, $v=\pi/2$, $b=v-\varphi$ and $T=\pi-\varphi$.
Let $f:[0,\pi]\to\R$ be absolutely continuous, with $f(v)=f(\pi)=0$.
Define, almost everywhere on the indicated intervals,
\begin{align*}
 r_1(t)&=-\tan t\,f(t)-f'(t), &&0<t<\varphi,\\
 r_2(t)&=f(t+v)-f'(t), &&\varphi<t<b,\\
 r_3(t)&=\frac{f(T)}{\sin(T-t)}-\cot(T-t)f(t)-f'(t), &&b<t<v,\\
 r_4(t)&=\cot t\,f(t)-f'(t), &&v<t<\pi.
\end{align*}
Suppose these residuals are square integrable, and let
$R^2=\sum_{j=1}^4\norm{r_j}_{L^2}^2$. Then
\begin{equation}\label{eq:stability-residual}
  \norm f_\infty\le\sqrt2\sec\varphi\,R.
\end{equation}
The constant is sharp for this pinned residual-space inequality.
\end{lemma}

\begin{proof}
The last residual equation and $f(v)=0$ give, for $v\le t<\pi$,
\[
 f(t)=-\sin t\int_v^t\frac{r_4(u)}{\sin u}\dd u.
\]
Cauchy--Schwarz bounds its square by
$-\sin t\cos t\,\norm{r_4}_{L^2}^2$. This tends to zero at $\pi$, so the
endpoint singularity is harmless. On $[b,v]$ another integrating factor gives
\[
 f(t)=-\frac{f(T)\cos t}{\cos\varphi}
       +\sin(T-t)\int_t^v\frac{r_3(u)}{\sin(T-u)}\dd u.
\]
On $[\varphi,b]$, substitute these identities into
$f(t)=f(b)-\int_t^b f(u+v)\dd u+\int_t^b r_2(u)\dd u$. Interchanging the
integrals yields
\begin{equation}\label{eq:stability-green-middle}
 f(t)=\int_t^b r_2(u)\dd u
     +\int_b^v\frac{r_3(u)}{\sin(T-u)}\dd u
     +\int_v^T k_t(u)r_4(u)\dd u,
\end{equation}
where
\[
 k_t(u)=
 \begin{cases}
  (\sec\varphi-\sin t)/\sin u,&v\le u\le v+t,\\
  (\sec\varphi+\cos u)/\sin u,&v+t\le u\le T.
 \end{cases}
\]
All denominators in this calculation are bounded away from zero. Finally,
on $[0,\varphi]$,
\[
 f(t)=\cos t\left(\frac{f(\varphi)}{\cos\varphi}
                    +\int_t^\varphi\frac{r_1(u)}{\cos u}\dd u\right).
\]

Each displayed formula expresses evaluation at $t$ as a linear functional
of $(r_1,r_2,r_3,r_4)$ in the direct sum of their four $L^2$ spaces. Its squared
norm is
\begin{equation}\label{eq:stability-diagonal}
 D(t)=
 \begin{cases}
 \cos^2t\,(2\sec^2\varphi-\tan t),&0\le t\le\varphi,\\
 \cos t\,(2\sec\varphi-\sin t),&\varphi\le t\le b,\\
 \sin t\cos t+2\tan\varphi\cos^2t,&b\le t\le v,\\
 -\sin t\cos t,&v\le t\le\pi.
 \end{cases}
\end{equation}
Here is the middle-interval computation, which also determines the first
interval. With $a=\sec\varphi$, the squared norm in
\eqref{eq:stability-green-middle} is
\[
 (b-t)+\tan\varphi+(a-\sin t)^2\tan t
       +\int_t^b(a-\sin w)^2\sec^2w\dd w.
\]
An antiderivative of the last integrand is
$(a^2+1)\tan w-2a\sec w-w$. Substituting
$\tan b=\cot\varphi$ and $\sec b=\csc\varphi$ gives the second line of
\eqref{eq:stability-diagonal}; in particular
$D(\varphi)=2-\sin\varphi\cos\varphi$. The first line follows by multiplying
this norm by $\cos^2t\sec^2\varphi$ and adding
$\cos^2t\int_t^\varphi\sec^2u\dd u$. On $[b,v]$, the two squared coefficient
norms are $\tan\varphi\cos^2t$ and
$\sin^2(T-t)(\tan\varphi-\cot(T-t))$, giving the third line.

The adjacent formulas agree at their endpoints. The first is at most
$2\sec^2\varphi$. The derivative of the second is
$-1-2\sin t(\sec\varphi-\sin t)\le-1$; the derivative of the third is
$\cos(2t)-4\tan\varphi\sin t\cos t\le0$, since $t\ge b>\pi/4$.
The last is at most $1/2$. Hence $\max D=D(0)=2\sec^2\varphi$, and
$|f(t)|^2\le D(t)R^2$ proves \eqref{eq:stability-residual}.

For sharpness, take the residual tuple to be the coefficient tuple for
evaluation at zero in these formulas. Its squared norm and the reconstructed
value $f(0)$ both equal $D(0)$. The coefficients are bounded and piecewise
smooth, with $r_4=0$ on $(T,\pi)$; reconstruction gives an absolutely
continuous $f$ with the prescribed endpoint values. Equality is therefore
attained. This does not assert sharpness within feasible cap perturbations,
or after minimizing over translations.
\end{proof}

\begin{proof}[Proof of \cref{thm:stability-cap}]
Let $\xi_0=\xi_G$ and $\xi_1=\xi$. Write
$\mathcal E=\frac12\sum_{j=1}^6\int(\varrho_j(\xi_1)-\varrho_j(\xi_0))^2$,
where the six terms are those of $\cS,\cR,\cL$. The affine decomposition
in \cref{fact:Q}\ref{fact:Q:decomp} and \eqref{eq:gap} give
\[
 \cQ((1-\lambda)\xi_0+\lambda\xi_1)
 =(1-\lambda)M+\lambda\cQ(\xi_1)+\lambda(1-\lambda)\mathcal E.
\]
The left side is at most $M$. Dividing by $\lambda>0$ and letting
$\lambda\downarrow0$ yields $M-\cQ(\xi_1)\ge\mathcal E$.

Put $f(t)=\Delta(t)-s\cos t$. Then $f(v)=f(\pi)=0$ by the top normalization
and the choice of $s$. Support functions are Lipschitz. The tangent-displacement
formulas used in \cref{lem:tangent} identify the differences in the four
cap terms with $r_1,\ldots,r_4$ of \cref{lem:stability-residual}. Translation
cancels from these differences. The Mamikon identities supply their square
integrability, so
\[
 \norm f_\infty^2
 \le 2\sec^2\varphi\sum_{j=1}^4\norm{r_j}_{L^2}^2
 \le4\sec^2\varphi\,\mathcal E
 \le4\sec^2\varphi\,(M-\cQ(\xi)).
\]
The two auxiliary-body squares were only discarded as nonnegative; no
strict-concavity assertion for all triple variables is needed.

On the lower semicircle the cap supports are determined by the two bottom
endpoints, so the same bound controls the full support functions. The
support-function formula for $\dH$ gives \eqref{eq:stability-Q}. Apply it to
the canonical triple and use $\cA(K)\le\cQ(\xi_K)$ to obtain
\eqref{eq:stability-area}. Finally
$\cos\varphi\ge1-\varphi^2/2\ge1249/1250$ implies
$2\sec\varphi\le2500/1249<1001/500$.
\end{proof}

\subsection{Qualitative stability without injectivity}

For nonempty compact sets, write
$d_{\rm rig}(S,G)=\inf_{\alpha,z}\dH(S,R_\alpha G+z)$, where here $\dH$
denotes the ordinary Hausdorff distance, not a support-function formula for
nonconvex sets.

\begin{theorem}[qualitative sofa stability]\label{thm:stability-global}
For every $\eta>0$ there is $\varepsilon>0$ such that every moving sofa $S$
with $\abs S\ge M-\varepsilon$ satisfies $d_{\rm rig}(S,\Gs)<\eta$.
Consequently any sequence with $\abs{S_n}\to M$ has rigid images converging
to $\Gs$ both in Hausdorff distance and in symmetric-difference area.
\end{theorem}

\begin{proof}
First we obtain a uniform bound on large sofas modulo translation. If a
compact connected set $S\subset\R\times[0,1]$ lies in its supporting hallway
at $t\in(0,\pi/2)$, its inner corner has height at most one. Indeed, in the
wall coordinates $a(p),b(p)$ relative to that corner, hallway membership says
$a,b\le1$ and $\max(a,b)\ge0$; both upper values are attained. If the corner
were above height one, the sets $\{a\ge0\}$ and $\{b\ge0\}$ would be disjoint,
nonempty, closed subsets covering $S$, contrary to connectedness.

By \cref{fact:angle}, sofas of area at least $11/5$ may be normalized into
the horizontal strip with top support one and rotation angle
$\omega\in[\omega_0,\pi/2]$. Since $\omega_0>\pi/4$, the preceding observation
applies at $\pi/4$. Writing $h_1=h_S(\pi/4)$ and $h_2=h_S(3\pi/4)$, it gives
$h_1+h_2\le2+\sqrt2$. The outer walls and the floor bound the abscissae by
$-\sqrt2h_2\le x\le\sqrt2h_1$. Center this interval by a horizontal translation.
Every such sofa then lies in
\[
 B=[-(1+\sqrt2),1+\sqrt2]\times[0,1].
\]
This preserves top support one, all supporting-hallway containments, and the
endpoint width inequality $h_S(\omega)+h_S(\omega+\pi)\le1$.

Suppose now $\abs{S_n}\to M$ but $d_{\rm rig}(S_n,\Gs)\ge\eta>0$. Normalize
as above. Compactness of the space of nonempty compact subsets of $B$ gives
a subsequence $S_n\to S$, with $S$ nonempty, compact and connected; also
$\omega_n\to\omega\in[\omega_0,\pi/2]$ after a further subsequence.
This is compact-hyperspace selection, not a claim that the sets are convex.
It follows, for example, by uniform subsequential convergence of their
bounded, $1$-Lipschitz distance functions on $B$.

Supports converge uniformly, even for nonconvex compact sets. They are
uniformly Lipschitz in angle because the sets lie in $B$. For each
$s\in[0,1]$ put $t_n=s\omega_n$ and $t=s\omega$. Passing to the limit in the
closed wall inequalities $a_n\le1$, $b_n\le1$, $\max(a_n,b_n)\ge0$ gives
$S\subset L_S(t)$. Top normalization and the endpoint width bound also pass
to the limit. Thus
\[
 \Phi_s(p)=R_{-s\omega}\bigl(p-\xx_S(s\omega)\bigr)
\]
is a continuous movement of $S$: the start lies in $H_L$, and at the end the
first wall coordinate lies in $[0,1]$ by the width bound and the second is at
most one. This constructs a motion of the limit rather than assuming
compactness of the original motion paths.

Area is upper semicontinuous here. For every $r>0$, eventually $S_n$ lies in
the closed $r$-neighborhood of $S$; continuity of measure from above as
$r\downarrow0$ gives $\abs S\ge\limsup_n\abs{S_n}=M$. Optimality gives
$\abs S=M$, and \cref{thm:main} makes $S$ a rigid image of $\Gs$.
Then $d_{\rm rig}(S_n,\Gs)\le\dH(S_n,S)\to0$, a contradiction.

Finally choose rigid images with $S_n\to\Gs$. Their excess outside $\Gs$ is
bounded by the area of a shrinking outer neighborhood of $\Gs$ minus $\Gs$,
hence tends to zero. The identity
$\abs{\Gs\setminus S_n}=M-\abs{S_n}+\abs{S_n\setminus\Gs}$ gives the same
conclusion for the missing area. No continuity of nonconvex area is assumed.
\end{proof}

\begin{remark}
The quantitative cap estimate does not supply an injectivity hypothesis or a
right-angle reduction for arbitrary near-maximizers. Accordingly
\cref{thm:stability-global} asserts no uniform power-law rate. The results of
this proposed section have not yet been included in the Lean formalization.
\end{remark}
